% from ega4-18-fr.tex
\oldpage[IV-4]{152}

  \item[(ii)] On ignore si les fibres formelles de $A$ sont géométriquement connexes (voire géométriquement irréductibles) lorsque $A$ est un anneau local hensélien quelconque. Il suffirait que, pour tout anneau local noethérien, hensélien et \emph{intègre} $A$, $\Spec(\widehat A)$ soit irréductible. On ignore s'il en est toujours ainsi.

  \item[(iii)] On ignore si lorsque $A$ est un anneau local excellent, son hensélisé strict ${}^{hs}\!A$ est excellent. On peut voir que, pour élucider cette question, on peut se ramener au cas où $A=k[[T_1,\ldots,T_n]]$, $k$ étant un corps de caractéristique $p>0$. Compte tenu de (18.8.17), la question est de savoir si les fibres formelles de ${}^{hs}\!A$ sont géométriquement régulières. La réponse est affirmative lorsque $n=1$, mais n'est pas connue pour $n=2$. On peut montrer que la réponse est affirmative chaque fois que, pour tout schéma $Y_1$ fini sur $Y=\Spec(\widehat A)$, on peut résoudre les singularités de $Y_1$ (7.9.1).
\end{enumerate}
\end{remarks}

L'assertion de connexité faite dans (18.9.1) se généralise de la façon suivante:

\begin{theorem}[18.9.7]
\label{IV.18.9.7}
Soient $A$, $B$ deux anneaux locaux noethériens, $\rho:A\to B$ un homomorphisme local, $f:\Spec(B)=X\to\Spec(A)=Y$ le morphisme correspondant. Supposons vérifiées les conditions suivantes:
\begin{enumerate}
  \item[(i)] $f$ est plat et toutes les fibres $X_y=f^{-1}(y)$ ($y\in Y$) sont géométriquement réduites (4.6.2).
  \item[(ii)] Ou bien l'anneau $A$ est géométriquement unibranche (0, 23.2.1), ou bien l'anneau $A$ est unibranche (0, 23.2.1) et $k(B)$ est une extension primaire (4.3.1) de $k(A)$.
\end{enumerate}
Alors, si $\eta$ est le point générique de $Y$, $X_\eta$ est connexe, et pour toute partie fermée $T$ de $X$ telle que $f(T)$ soit rare dans $Y$, $X-T$ est connexe.
\end{theorem}

Les deux conclusions énoncées sont en fait équivalentes; en effet, l'hypothèse que $f(T)$ est rare entraîne que $f(T)$ ne contient pas $\eta$; d'autre part, comme $\{\eta\}$ est proconstructible (I, 9.6) et que $f$ est plat, $X_\eta$ est dense dans $X$ (2.3.10), et par suite si $X_\eta$ est connexe, il en est de même de $X-T$ qui le contient. Inversement, lorsque $U$ parcourt l'ensemble des voisinages ouverts affines de $\eta$ dans $Y$, $X_\eta$ est limite projective des sous-préschémas induits sur les ouverts $f^{-1}(U)$ de $X$ (8.1.2, $a)$); si l'on suppose que $X-T$ est connexe pour toute partie fermée $T$ de $X$ telle que $f(T)$ soit rare dans $Y$, les $f^{-1}(U)$ sont connexes, donc il en est de même de $X_\eta$ (8.4.1).

Montrons d'abord qu'on peut se limiter au cas où $A$ est intègre et $Y=\Spec(A)$ géométriquement unibranche (6.15.1) (ce qui implique que $A$ est géométriquement unibranche, mais n'est pas équivalent à cette condition (6.15.2)). Il est clair qu'on peut d'abord supposer $Y$ réduit, en considérant le morphisme $X\times_Y Y_{\mathrm{red}}\to Y_{\mathrm{red}}$, déduit de $f$ par changement de base, qui est plat, et a mêmes fibres que $f$. On peut donc supposer $A$ intègre et unibranche; si $K$ est le corps des fractions de $A$, il existe alors une sous-$A$-algèbre finie $A''$ de $K$ telle que, si $A'$ est la clôture intégrale de $A$, le morphisme $\Spec(A')\to\Spec(A'')$ soit radiciel (0, 23.2.5); on en conclut (6.15.5) que $\Spec(A'')$ est géométriquement unibranche. Comme $A$ est unibranche, $A'$ est un anneau local, donc il en est de même de $A''$; montrons que l'anneau $B''=B\otimes_A A''$ est aussi un anneau local. En effet, $B''$ est une $B$-algèbre finie, donc un anneau semi-local noethérien; si $k''$ est le corps résiduel de $A''$, les idéaux maximaux de $B''$ sont les points de $\Spec(B'')$ au-dessus du point fermé de $\Spec(B)$, donc les points de $\Spec(k(B)\otimes_{k(A)}k'')$ (I, 3.4.9)

